% id: 140-17
% book: 베이직쎈 미적분1 · 09 정적분의 활용
% section: 개념 57 두 곡선 사이의 넓이
% passage: 다음 그림에서 색칠한 부분의 넓이를 구하시오.
% figure: tikz:fig-140-17
% answer: 
% answer_source: 
% variant_level: 0 (원본 전사)
% 이 파일은 build.mjs 가 items.json 에서 만든다. 고칠 때는 items.json 을 고친다.
\smash{\raisebox{1.5mm}{\dmkichul{베이직쎈 미적분1 140쪽 17번}}}
\noindent\begin{minipage}[t]{0.58\linewidth}\vspace{0pt}\raggedright \end{minipage}\hfill\begin{minipage}[t]{0.4\linewidth}\vspace{0pt}\centering \resizebox{\ifdim\width>\linewidth\linewidth\else\width\fi}{!}{% 140-17(140쪽) — 곡선 $y=x^2-2x$ 와 직선 $y=x$ 로 둘러싸인 색칠한 부분($x=0$ 에서 $x=3$ 까지). 스캔 화질이 낮아 TikZ 로 다시 그림.
% 원문에 있는 것만 옮긴다: 포물선 · 직선 · 색칠 · $x=3$ 의 안내 점선 · 지시선 · 라벨 y=x^2-2x, y=x, O, 2, 3, x, y. 원문에 없는 눈금·넓이 값은 적지 않는다.
\begin{tikzpicture}[x=0.32cm, y=0.26cm, line width=0.6pt]
  \fill[gray!25] plot[domain=0:3, samples=80, smooth] (\x, {(\x)}) -- plot[domain=3:0, samples=80, smooth] (\x, {(\x)*(\x)-2*(\x)}) -- cycle;
  \draw[->, >=stealth, line width=0.7pt] (-2.3,0) -- (4.2,0) node[below=1pt, font=\small] {$x$};
  \draw[->, >=stealth, line width=0.7pt] (0,-1.9) -- (0,6.3) node[left=1pt, font=\small] {$y$};
  \draw[densely dashed, line width=0.4pt] (3,0) -- (3,3);
  \draw[domain=-1.65:3.15, samples=100, smooth] plot (\x, {(\x)*(\x)-2*(\x)});
  \draw (-1.9,-1.9) -- (3.7,3.7);
  \draw[->, >=stealth, line width=0.4pt] (2.5,4.9) .. controls (3.0,4.5) and (3.0,3.8) .. (2.95,2.95);
  \node[anchor=west, inner sep=1pt, font=\small] at (1.1,5.6) {$y=x^2-2x$};
  \node[anchor=west, inner sep=1pt, font=\small] at (3.6,3.75) {$y=x$};
  \node[below=1pt, font=\small] at (2,0) {$2$};
  \node[below=1pt, font=\small] at (3,0) {$3$};
  \node[below left=-2pt and -2pt, font=\small] at (0,0) {$\pt{O}$};
\end{tikzpicture}
}\end{minipage}\par
\dmhint{$\displaystyle\int_{0}^{\text{\scriptsize\setlength{\fboxsep}{1.5pt}\framebox[1.5em]{\rule{0pt}{1.6ex}}}} \{\blank[1.1em]-(x^2-2x)\}\mkern3mu dx$\\ $\displaystyle =\int_{0}^{\text{\scriptsize\setlength{\fboxsep}{1.5pt}\framebox[1.5em]{\rule{0pt}{1.6ex}}}} (-x^2+\blank[1.1em])\mkern3mu dx$\\ $\displaystyle =\Big[-\dfrac{1}{3}x^3+\blank[1.1em]x^2\Big]_{0}^{\text{\scriptsize\setlength{\fboxsep}{1.5pt}\framebox[1.5em]{\rule{0pt}{1.6ex}}}}$\\ $=\blank$}
